% !PRINTED 478
\weberpage{478}
If one has any system of non-vanishing functions $X_1, X_2, X_3$, one can therefore always form such resolvents. By the form problem of the group $G_{168}$, then, the equation of degree 7 with the group $P_{168}$ is solved at the same time.

In order to reduce the solution of these equations of degree 7 to the special form problem as it occurs in the elliptic functions (with $f = 0$), an accessory equation of degree 4 must then still be solved, just as, in order to reduce the general equation of degree 5 to the icosahedral equation, we found an accessory square root necessary.

Everything now depends, then, on determining the functions $X_1, X_2, X_3$ of the roots $\lambda$ in accordance with these requirements. To make this possible, we must make use of some simple theorems from the general theory of invariants, which we shall derive in the following section, so far as they come into consideration for our task.

\begin{center}
§. 129.

\gesperrt{Contragredient Groups.}
\end{center}

We have already explained the concept of the contragredient transformation in §.\,37. Namely, if

\weq{1}{A = \begin{pmatrix} a_1, & b_1, & c_1 \\ a_2, & b_2, & c_2 \\ a_3, & b_3, & c_3 \end{pmatrix}, \quad A_1 = \begin{pmatrix} a_1, & a_2, & a_3 \\ b_1, & b_2, & b_3 \\ c_1, & c_2, & c_3 \end{pmatrix}}

are two substitutions transposed to each other, and $x_1, x_2, x_3$ and $\xi_1, \xi_2, \xi_3$ are two sequences of variables which, by the substitutions

\weq{2}{(y) = A(x), \quad \xi = A_1(\eta)}

are transformed into two new sequences of variables $y_1, y_2, y_3$ and $\eta_1, \eta_2, \eta_3$, then we have called these two sequences of variables, and likewise their transformations, \gesperrt{contragredient}.

By the substitution $y = A(x)$ every function $\Phi(x_1, x_2, x_3)$ of the $(x)$ is transformed into a function of the $(y)$, and the formation of the derivatives yields:

\weq{3}{\frac{\partial \Phi}{\partial x_1} = a_1 \frac{\partial \Phi}{\partial y_1} + a_2 \frac{\partial \Phi}{\partial y_2} + a_3 \frac{\partial \Phi}{\partial y_3}, \cdots}
%CONT>>
